% 卷二完整题解；由 main.tex 在全文末尾载入。
\begingroup
\setlength{\emergencystretch}{3em}

先完成题目，再来核对。下面按正文顺序整理随堂练习与章末题；两套编号分开，例如“随堂练习 1.1”与“第一章章末第 1 题”不是同一道题。
操作题给出步骤、预期结果和判断依据；依赖设备、软件版本或实验数据的数值，以实际记录为准。
本章验收用于现场考核，不另列一份可以照抄的验收答案。

\section*{第一章：ROS2 基础}
\phantomsection
\addcontentsline{toc}{section}{第一章题解：ROS2 基础}
\subsubsection*{随堂练习}

\paragraph*{随堂练习 \ref{ex:v2-ros-setup}}
\begin{solution}
talker 属于 \texttt{demo\_nodes\_cpp}。加载当前发行版环境后，可执行：
\begin{tbash}
ros2 pkg executables demo_nodes_cpp
ros2 run --help
ros2 run demo_nodes_cpp talker
\end{tbash}
前两条分别确认包中的可执行程序与启动语法，第三条启动发布节点。按 Ctrl+C 结束。
\end{solution}

\paragraph*{随堂练习 \ref{ex:v2-ros-graph}}
\begin{solution}
在只启动一个示例 talker 的环境中，\texttt{/chatter} 类型为 \texttt{std\_msgs/msg/String}，发布者数量通常为 1。
\begin{tbash}
ros2 topic info /chatter
ros2 node list
\end{tbash}
停止 talker 后等待发现信息更新，发布者数量应变为 0，或话题不再列出；如果 echo 等订阅者仍在运行，话题可能仍可见。节点是功能实体，话题是通信通道，不能用“还能查到话题”判断发布节点是否存活。
\end{solution}


\paragraph*{随堂练习 \ref{ex:v2-ros-qos}}
\begin{solution}
可靠发布者能够满足 best\_effort 订阅者的要求，所以可以通信。发布端和订阅端按“提供的策略是否满足要求”匹配，不要求每项值都相等。分别运行：
\begin{tbash}
# 终端一
ros2 topic pub /qos_demo std_msgs/msg/String "{data: hello}" --rate 1 --qos-reliability reliable
# 终端二
ros2 topic echo /qos_demo std_msgs/msg/String --qos-reliability best_effort --once
\end{tbash}
应收到 hello。这里只改变可靠性，其他策略保持兼容；不能由这个实验推断任何 QoS 组合都能通信。
\end{solution}

\paragraph*{随堂练习 \ref{ex:v2-ros-message}}
\begin{solution}
在 \texttt{my\_msgs/msg/My.msg} 末尾加入 \texttt{float64 distance}，保存后回到工作空间根目录：
\begin{tbash}
cd ~/ros2_ws
colcon build --packages-select my_msgs
source install/setup.bash
ros2 interface show my_msgs/msg/My
\end{tbash}
应看到新增字段。若节点要访问该字段，再重新构建依赖它的包，例如执行 \texttt{colcon build --packages-up-to my\_package}。单位约定为米，应同时写在接口注释或文档里。
\end{solution}

\paragraph*{随堂练习 \ref{ex:v2-ros-timer}}
\begin{solution}
将构造函数中的间隔改为 \texttt{std::chrono::milliseconds(500)}，重新构建、加载环境并重启节点。在另一个已加载环境的终端执行：
\begin{tbash}
ros2 topic hz /timer_
\end{tbash}
理论频率为 2 Hz，调度和观测窗口会造成小幅波动。计数器只表示回调执行次数；时间戳表示特定时钟下的事件时间。间隔改变后，一次计数不再等于一秒，不能用计数器代替时间戳。
\end{solution}

\paragraph*{随堂练习 \ref{ex:v2-ros-service}}
\begin{solution}
服务端运行时查询：
\begin{tbash}
ros2 service list -t
ros2 interface show example_interfaces/srv/AddTwoInts
\end{tbash}
请求包含两个 \texttt{int64} 字段 a、b，响应包含 \texttt{int64 sum}；接口中的三条短横线分隔请求和响应。服务对应一次请求与一次响应，话题用于按发布端节奏持续传递数据，两者用途不同。
\end{solution}


\paragraph*{随堂练习 \ref{ex:v2-ros-reply}}
\begin{solution}
receive 与 send 都使用 \texttt{std\_msgs/msg/String}。先保持 \texttt{ros2 topic echo /send} 运行，再在另一终端发送：
\begin{tbash}
ros2 topic pub --once /receive std_msgs/msg/String "{data: armor_found}"
\end{tbash}
应收到 \texttt{I can hear:armor\_found}。停止节点后，原命令可能等待订阅者；加 \nolinkurl{--wait-matching-subscriptions} 并设为 0，可以在没有订阅者时也发送一次，但不会收到回复。CLI 打印 publishing 只说明发布端执行了发送，处理是否发生要由回调日志、响应消息等证据确认。
\end{solution}

\paragraph*{随堂练习 \ref{ex:v2-ros-domain}}
\begin{solution}
域不同的两个节点不能按普通 DDS 发现互相通信；统一为 42 后应恢复收发。两端都先停止节点，再设置相同的 \texttt{ROS\_DOMAIN\_ID} 和发现范围，执行 \texttt{ros2 daemon stop} 并重新启动。恢复默认实验环境时停止节点、unset \texttt{ROS\_DOMAIN\_ID}、\texttt{ROS\_LOCALHOST\_ONLY}、\nolinkurl{ROS_AUTOMATIC_DISCOVERY_RANGE} 和 \texttt{ROS\_STATIC\_PEERS}，停止 daemon 后重新运行；持久写入启动文件的设置也需删除或更改。域隔离实验须排除同名遗留发布者。
\end{solution}

\paragraph*{随堂练习 \ref{ex:v2-ros-ament}}
\begin{solution}
启动系统通常从安装索引查资源。检查 CMake 是否安装 launch/config 目录、是否重新构建，以及当前环境是否指向该工作空间。手动版本应包含：
\begin{tcode}
install(DIRECTORY launch config DESTINATION share/${PROJECT_NAME})
\end{tcode}
间隔函数改为 \texttt{milliseconds >= 10 \&\& milliseconds <= 500}；合法边界测试为 10 和 500，非法测试为 9 和 501。\Needspace{11\baselineskip}
保存后执行：
\begin{tbash}
cd ~/ros2_ws
colcon build --packages-select training_tools --cmake-args -DBUILD_TESTING=ON
source install/setup.bash
colcon test --packages-select training_tools
colcon test-result --verbose
ros2 pkg prefix training_tools
\end{tbash}
预期边界测试全部通过，包前缀指向当前工作空间的安装目录。
\end{solution}

\paragraph*{随堂练习 \ref{ex:v2-ros-launch}}
\begin{solution}
在 team 文件中将 observer 的 cmd 改成下面的列表，其余配置保持不变。嵌套列表把三个替换片段拼成同一个话题参数：
\begin{tpython}
cmd=['ros2', 'topic', 'echo',
     ['/', LaunchConfiguration('robot'), '/health'],
     'std_msgs/msg/String', '--once'],
\end{tpython}
重新构建、加载环境后运行：
\begin{tbash}
ros2 launch training_tools team.launch.py robot:=blue
\end{tbash}
应看到一条 team ready 消息及观测进程退出提示，heartbeat 继续运行。结束后再次启动：
\begin{tbash}
ros2 launch training_tools team.launch.py robot:=blue watch:=false
\end{tbash}
此时不启动 echo 进程。用另一个终端查询 \texttt{ros2 node list}、\texttt{ros2 topic info /blue/health}，应确认 heartbeat 发布者存在；查询命令本身可能产生临时观测节点，不能仅凭节点总数判断。
\end{solution}

\paragraph*{随堂练习 \ref{ex:v2-ros-log}}
\begin{solution}
定时器仍为 200 毫秒，所以理论消息频率是 5 Hz；alive 限流改成 1000 毫秒后约为 1 Hz。修改源码、重新构建、加载环境并重启后，用 \texttt{ros2 topic hz /status} 测量前者，按日志时间戳检查后者。调度、观测窗口和所用时钟会影响实测间隔。WARN 阈值仅隐藏 INFO/DEBUG 日志，发布逻辑仍运行，不会因此停止消息。
\end{solution}

\paragraph*{随堂练习 \ref{ex:v2-ros-ide}}
\begin{solution}
IDE 默认运行独立构建目录中的目标，\texttt{ros2 run} 根据已加载环境运行安装目录中的目标。仅执行 IDE Build 时，安装目录可能仍是旧程序。先去掉运行配置中的 \texttt{-p text:=from\_ide} 等文本覆盖，才能观察修改后的默认值；否则参数会掩盖源码默认值的变化。停止旧进程，然后执行：
\begin{tbash}
cd ~/ros2_ws
source /opt/ros/jazzy/setup.bash
colcon build --packages-select training_tools
source install/setup.bash
ros2 pkg prefix training_tools
ros2 run training_tools heartbeat
\end{tbash}
确认包路径属于当前工作空间，观察 /status 应得到 from\_debug。若仍是旧值，检查是否有其他发布者、参数覆盖、同名覆盖包，以及实际改动的源码路径。
\end{solution}

\subsubsection*{章末习题}

\paragraph*{第一章章末第 \ref{prob:v2-ros-1} 题}
\begin{solution}
包为 1，两个程序各启动一次时进程与节点各为 2；同一进程也可以包含多个节点。
\end{solution}

\paragraph*{第一章章末第 \ref{prob:v2-ros-2} 题}
\begin{solution}
加载 install/setup.bash；build 是中间构建结果，log 是日志。
\end{solution}

\paragraph*{第一章章末第 \ref{prob:v2-ros-3} 题}
\begin{solution}
不是，它表示历史队列深度；周期发布通常由定时器及回调执行情况决定。
\end{solution}

\paragraph*{第一章章末第 \ref{prob:v2-ros-4} 题}
\begin{solution}
类型为 \texttt{my\_msgs::msg::My}，包含头文件 \texttt{my\_msgs/msg/my.hpp}。
\end{solution}

\paragraph*{第一章章末第 \ref{prob:v2-ros-5} 题}
\begin{solution}
核对消息类型、命名空间、域标识、QoS 兼容性和终端环境；再确认节点存活、发布逻辑与执行器是否运行。
\end{solution}

\paragraph*{第一章章末第 \ref{prob:v2-ros-6} 题}
\begin{solution}
响应 sum 应为 5；这是一问一答服务，不是持续推送的话题。
\end{solution}

\paragraph*{第一章章末第 \ref{prob:v2-ros-7} 题}
\begin{solution}
环境变量由新进程继承，已启动的节点不会读取另一个终端后来修改的环境。应停止旧节点、重新设置通信范围、停止 daemon 后重启。域编号用于分组，本机范围用于限制通信接口；两者都没有提供身份认证，也不能阻止程序自建的其他网络连接。
\end{solution}

\paragraph*{第一章章末第 \ref{prob:v2-ros-8} 题}
\begin{solution}
完整话题名为 \texttt{/blue/health}。检查 launch/config 的安装规则、构建结果、环境加载顺序及 \texttt{ros2 pkg prefix training\_tools} 返回路径；再检查 get\_package\_share\_directory 拼出的文件名。源目录存在文件不代表安装目录已有文件。YAML 节点键还需匹配实际节点名，或使用本例的通配键。
\end{solution}

\paragraph*{第一章章末第 \ref{prob:v2-ros-9} 题}
\begin{solution}
消息理论频率为 5 Hz，限流日志约为 0.5 Hz。日志限流不改变定时器回调。FATAL 只记录严重级别，不自动终止程序；程序应明确执行故障处置和退出逻辑。
\end{solution}

\paragraph*{第一章章末第 \ref{prob:v2-ros-10} 题}
\begin{solution}
先确认修改的源码属于哪个包、IDE 实际运行哪个构建目标，再确认新目标是否已安装，最后检查当前终端加载的前缀及启动参数。设置 CMAKE\_INSTALL\_PREFIX 不等于执行安装。停止旧节点后可用以下命令统一结果：
\begin{tbash}
cd ~/ros2_ws
source /opt/ros/jazzy/setup.bash
colcon build --packages-select training_tools
source install/setup.bash
ros2 pkg prefix training_tools
ros2 run training_tools heartbeat
\end{tbash}
若有多个覆盖工作空间，使用只加载系统环境和目标工作空间的新终端复核；同时去掉覆盖默认文本的参数，检查话题上是否还有旧发布者。
\end{solution}

\section*{第二章：多机通信方法}
\phantomsection
\addcontentsline{toc}{section}{第二章题解：多机通信方法}
\subsubsection*{随堂练习}

\paragraph*{随堂练习 \ref{ex:v2-comm-boundary}}
\begin{solution}
可让相机处理进程作为 Server、记录进程作为 Client，使用本机 TCP；协议约定长度、字段类型/编码、序号、时间、单位与坐标系。跨机器还约定地址与访问控制。设备由一个明确进程读取后分发完整帧，否则两个读取者可能分别拿走部分字节，双方都无法组帧。
\end{solution}

\paragraph*{随堂练习 \ref{ex:v2-comm-roles}}
\begin{solution}
Server 执行 bind、listen、accept，Client 执行 connect；随后双方都能收发。把 ping 改成 pong 仍输出 pong。改为两字节后还要关闭发送方向，否则只是暂时没收满，会等待到超时：
在 Python 源码中修改为下面两行：
\begin{tpython}
client.sendall(b"po")
client.shutdown(socket.SHUT_WR)
\end{tpython}
接收累计到两字节后，下一次 recv 得到空字节，抛出“连接提前关闭”；不能把 po 当完整四字节结果。
\end{solution}

\paragraph*{随堂练习 \ref{ex:v2-comm-fragments}}
\begin{solution}
把 \texttt{parser\_demo} 中的输入列表替换为：
\begin{tpython}
[b"\x02", b"A", b"B", b"\x03X", b"YZ"]
\end{tpython}
输出仍为 AB、XYZ。向 LengthParser.feed 输入零长度会抛出 ValueError 并清空缓存，TCP 服务端据此结束本连接。协议没有可靠帧头，后面某个载荷字节也可能像合法长度，跳过一个字节再猜会误组帧。
\end{solution}

\paragraph*{随堂练习 \ref{ex:v2-comm-latest}}
\begin{solution}
priority 改为 status、position 后，第一次 take 得到 status/ready，第二次得到 position/30，槽数随后为零；调整原来只检查位置的断言和输出。代次 1 的迟到清理应记录当前代次 2、返回 False，位置 42 保留；代次 2 清理返回 True，清空自身状态。实际共享状态的检查与修改要在同一保护范围内，日志的状态读取也应同步。运行后面的两轮 TCP，第二轮同样拒绝代次 1 清理并继续回显。
\end{solution}

\paragraph*{随堂练习 \ref{ex:v2-comm-events}}
\begin{solution}
把丢确认条件改为 \texttt{event[0] == 1}，不再限定 attempt 为 1。第 1 事件发送三次，实际只执行一次，每次确认都丢失；最终 RuntimeError 报告投递失败，程序不会继续处理第 2 事件。发送端失败并不证明接收端没执行，应通过约定的查询/确认处理不确定结果。原示例只丢第一次时，第 1 事件第二次得到确认，第 2 事件也完成。
\end{solution}

\paragraph*{随堂练习 \ref{ex:v2-comm-crc}}
\begin{solution}
不要重算损坏帧的 CRC。在 \texttt{parser\_demo} 构建 stream 处改为：
\begin{tpython}
damaged = bytearray(encode_crc(7))
damaged[1] ^= 1
stream = b"noise" + bad + damaged + encode_crc(9)
\end{tpython}
值列表应为 [9]，坏 CRC 计数为 2，最后 buffer 为空；同步修改原本期待 [7] 与一次坏 CRC 的断言。改完数据又重新算 CRC，得到的是另一条有效帧，不能测试拒收损坏载荷。
\end{solution}

\paragraph*{随堂练习 \ref{ex:v2-comm-units}}
\begin{solution}
\texttt{bridge\_demo} 中 struct.pack 的输入改为 (8, 2000)。普通模式输出 (8, 2.0)，ROS 的 /\texttt{range\_m} 输出 data: 2.0；换程序前停止旧节点，避免两发布者混在一起。换算在协议边界做一次，内部统一使用米。位置接口还需明确坐标轴/坐标系、时间来源、单位、序号与无效值；Float64 的类型本身不会帮你约定这些内容。
\end{solution}

\paragraph*{随堂练习 \ref{ex:v2-comm-isolation}}
\begin{solution}
A 最后一次输入在第 1 秒，第 4 秒年龄为 3。阈值 4 时不过期，应得到 position/11，expired 保持 0；相应修改演示中原本期待 None 和计数 1 的断言，并把过期提示改为“旧值仍被允许”。恢复阈值 2 时丢弃一次。B 始终消费 100--105；A 槽数最多 1，第 5 秒消费新值 15。比较是大于阈值，所以年龄恰好等于阈值时仍允许消费。
\end{solution}

\subsubsection*{章末习题}
\paragraph*{第二章章末第 \ref{prob:v2-comm-1} 题}
\begin{solution}
不应争抢设备读取。由唯一所有者校验组帧，再向其他进程分发完整消息；发送控制命令也按同一接口约定协调。
\end{solution}
\paragraph*{第二章章末第 \ref{prob:v2-comm-2} 题}
\begin{solution}
客户端 connect，服务端 accept 得到新连接对象，通过它收发；监听对象继续负责接入连接。
\end{solution}
\paragraph*{第二章章末第 \ref{prob:v2-comm-3} 题}
\begin{solution}
不能。两条连接的数据没有组成同一帧的约定，拼起来可能混合不同轮次或不同客户端的数据。丢弃旧半帧，新连接另建解析器。
\end{solution}
\paragraph*{第二章章末第 \ref{prob:v2-comm-4} 题}
\begin{solution}
校验失败说明候选帧不满足协议，拒收并按规则重同步；CRC 不验证身份，也不能替代鉴权。
\end{solution}
\paragraph*{第二章章末第 \ref{prob:v2-comm-5} 题}
\begin{solution}
填 3.0 m/s，还需约定速度向量的坐标系、采集时间与时钟来源。不能把接收时间冒充源采集时间，也不能重复换算。
\end{solution}
\paragraph*{第二章章末第 \ref{prob:v2-comm-6} 题}
\begin{solution}
检查之后、关闭之前可能已换连接；在锁外重新读取当前 socket，就会关掉新对象。比较代次和摘下本轮归属要一起受保护，关闭保存的旧对象，不再读取当前对象。槽位写入和清理也遵守同样规则。
\end{solution}
\paragraph*{第二章章末第 \ref{prob:v2-comm-7} 题}
\begin{solution}
只消费最新值 300；三次购买需有界事件队列，约定确认、重试、序号去重及满队列/失败处理，不能覆盖成一次。
\end{solution}
\paragraph*{第二章章末第 \ref{prob:v2-comm-8} 题}
\begin{solution}
原示例第 1 事件发送 2 次、执行 1 次、确认成功；第 2 事件发送/执行各 1 次。所有确认都丢时，第 1 事件发送 3 次、执行 1 次，发送端失败并停止，第二事件尚未执行。检查日志中实际执行与发送次数，失败不能推出对方未执行；需要查回结果或按协议处理不确定状态。
\end{solution}
\paragraph*{第二章章末第 \ref{prob:v2-comm-9} 题}
\begin{solution}
年龄为 3 秒，超过 2 秒阈值，应丢弃、计数并按约定使用新鲜备用源。不同时间原点和未同步时钟不能直接相减；接收年龄使用本机同一单调时钟，跨机采集时间另行约定同步与误差范围。
\end{solution}


\section*{第三章：SDR 与电磁波解析}
\phantomsection
\addcontentsline{toc}{section}{第三章题解：SDR 与电磁波解析}
\subsubsection*{随堂练习}

\paragraph*{随堂练习 3.1}
\begin{solution}
天线与 ADC 属于硬件采集环节，相关峰搜索与 CRC 校验属于数字处理。完整记录下来的 IQ 可用于重放滤波、同步、解调和帧校验；无法补回采样前已丢失的信号，也不能重现天线摆放、模拟饱和等硬件条件的全部影响。
\end{solution}

\paragraph*{随堂练习 3.2}
\begin{solution}
将 f 改为 -125000 后，相邻样本相位差为 $-\pi/4$，约 -0.785398 rad；原断言需相应改为比较负值。正、负频率在复平面上旋转方向相反，但理想样本幅值都为 1，所以只看幅值无法区分。
\end{solution}

\paragraph*{随堂练习 3.3}
\begin{solution}
在状态转移查询前处理 reset，使其不依赖当前状态：
\begin{tpython}
if event == "reset":
    state = "level1"
else:
    state = transitions.get((state, event), state)
\end{tpython}
可依次输入 ok1、reset、ok2，状态应为 level2、level1、level1。超时重试通常在当前阶段重新尝试；主动复位返回初始阶段，必要时还应清空阶段缓存。
\end{solution}

\paragraph*{随堂练习 3.4}
\begin{solution}
原比特 1010 变为 0101，因为乘以 -1 等价于相位旋转 $\pi$，实部符号全部相反。接收机需要建立符号定时、载波相位与帧起点，并用已知前导等方式消除歧义；仅按正负判决不足以组成可靠接收机。
\end{solution}

\paragraph*{随堂练习 3.5}
\begin{solution}
两路数据年龄均为 $10.0-9.0=1.0$ s，超过 0.2 s 有效期，selected 应为 None。验证时应将原断言改为 \texttt{assert selected is None}。继续沿用最后坐标会把历史位置伪装成当前状态，应报告无有效源或进入明确的回退状态。
\end{solution}

\paragraph*{随堂练习 3.6}
\begin{solution}
一个明确的保存约定可以是：采样率 8000 Hz，样本为交错 I/Q 的小端 float32 数对，复数样本数 32；若无文件头，数据区应为 $32\times2\times4=256$ 字节。缺采样率会误解时间与频率，缺数据类型会误解数值和步长，字节序错误会破坏数值，样本数错误可能把截断或尾部数据误当成有效采样。保存时还应写明 I/Q 顺序与是否含文件头。
\end{solution}

\subsubsection*{章末习题}

\paragraph*{第三章章末第 1 题}
\begin{solution}
不能。离线处理只能使用已采集的信息，软件不能恢复从未记录的频段。
\end{solution}

\paragraph*{第三章章末第 2 题}
\begin{solution}
幅值为 5；会丢失相位信息，不能完整表示复数基带信号。
\end{solution}

\paragraph*{第三章章末第 3 题}
\begin{solution}
不应。事件与当前阶段不匹配，应保持状态或按明确策略报错。
\end{solution}

\paragraph*{第三章章末第 4 题}
\begin{solution}
同步确定符号或帧参考，解调把样本映射为比特或软信息，CRC 检查候选帧的数据一致性。
\end{solution}

\paragraph*{第三章章末第 5 题}
\begin{solution}
不可，样本年龄 0.4 s 超过有效期，应按约定回退或报告无有效源。
\end{solution}

\paragraph*{第三章章末第 6 题}
\begin{solution}
判决比特翻转；幅值不变，需要载波相位参考、前导或其它消歧机制。
\end{solution}

\paragraph*{第三章章末第 7 题}
\begin{solution}
相位差为 π/4 rad，共 4 个周期；应记录复数类型、采样率、频率、样本数与保存格式。
\end{solution}

\section*{第四章：OpenCV 与图像处理基础}
\phantomsection
\addcontentsline{toc}{section}{第四章题解：OpenCV 与图像处理基础}
\subsubsection*{随堂练习}

\paragraph*{随堂练习 4.1}
\begin{solution}
把 add\_executable、target\_include\_directories 和 target\_link\_libraries 中的目标名全部改为 image\_demo，不能只改其中一行。
\begin{tbash}
cmake -S . -B build
cmake --build build --target image_demo
./build/image_demo
\end{tbash}
源码目录为当前目录，包含 main.cpp 和 CMakeLists.txt；构建目录为其下的 build。目标名称与 project 名称可以不同。
\end{solution}

\paragraph*{随堂练习 4.2}
\begin{solution}
把 b 定义为 \texttt{cv::Mat b = a(cv::Rect(0,0,1,1));}，仍通过 b 修改其第一个像素。a 的对应像素变为 9，事先 clone 得到的 c 仍为 1。ROI 创建的是共享原数据的表头，clone 分配独立数据，因此修改传播范围不同。
\end{solution}

\paragraph*{随堂练习 4.3}
\begin{solution}
仅把 imread 的路径改为不存在的文件，保留前面的图像生成步骤。程序应在 empty 检查处返回 2。在 Linux 终端运行程序后立即执行 \texttt{echo \$?}，应输出 2。空 Mat 没有有效像素数据，继续转换颜色或访问像素可能触发断言或错误。
\end{solution}

\paragraph*{随堂练习 4.4}
\begin{solution}
第二个像素 H=90，不属于任一红色区间，第一次修改后 countNonZero 为 1。再将第一个像素 S 设为 0，它低于饱和度下限 70，非零数变为 0。色调、饱和度和亮度条件需要同时满足。
\end{solution}

\paragraph*{随堂练习 4.5}
\begin{solution}
输入为 126、127、128，阈值为 127。反向二值阈值在输入小于或等于阈值时输出 255，大于阈值时输出 0，因此结果为 255、255、0。边界值 127 属于前一类。
\end{solution}

\paragraph*{随堂练习 4.6}
\begin{solution}
改用 \texttt{cv::blur(a, out, cv::Size(3,3));}。中心邻域中一个值为 255、八个为 0，均值为 $255/9\approx28.33$，存回 8 位图时中心约为 28；中值是排序后的第五个值，即 0。均值会让孤立点影响邻居，中值能够在该例中直接消除它。
\end{solution}

\paragraph*{随堂练习 4.7}
\begin{solution}
7×7 核大于 5×5 方块，腐蚀阶段没有任何位置能完全容纳结构元素，方块和孤立点都被去掉；后续膨胀不能恢复已经消失的前景，非零数为 0。核越大越可能删掉真正的小目标或细结构。
\end{solution}

\paragraph*{随堂练习 4.8}
\begin{solution}
先对 a 执行 Canny，再在边缘图上使用相同的外轮廓模式，打印轮廓数量、各轮廓面积与边缘非零像素数。可使用阈值 50、150；具体边界位置及数值以该版本实测为准。原实心图的非零像素数是 200，而几何轮廓面积是 171；Canny 输出的是边缘像素，既不是实心区域计数，也不保证与原轮廓几何面积相等。
\end{solution}

\paragraph*{随堂练习 4.9}
\begin{solution}
把缩放插值改为 \texttt{cv::INTER\_NEAREST}，保存为另一个文件与 small.png 对照。最近邻只复制源像素，边缘容易呈台阶；INTER\_AREA 在缩小时按覆盖面积综合像素，更适合彩色图像下采样。对于类别掩膜，通常用最近邻避免产生不存在的类别值。
\end{solution}

\paragraph*{随堂练习 4.10}
\begin{solution}
以平移矩阵 $\left[\begin{smallmatrix}1&0&5\\0&1&0\end{smallmatrix}\right]$ 调用 warpAffine 得到第二张图，分别计算两图 ORB 描述子，再用 Hamming 距离匹配。记录两图特征数、匹配数与匹配点坐标差；正确对应大多应满足横向约 5 pixel、纵向约 0 的位移。数量受裁剪、检测器与版本影响。重复纹理也可能有低描述子距离，因此还需几何一致性检查。
\end{solution}

\paragraph*{随堂练习 4.11}
\begin{solution}
保留原消息编码 bgr8，把 \texttt{toCvCopy(msg,"bgr8")} 的目标编码改为 mono8，并打印 \texttt{out->image.channels()}，应为 1。转换函数执行颜色到灰度的变换；仅把消息编码标签改成 mono8，而仍放三通道数据，会使声明与数据不符。编码还描述颜色顺序、元素深度和语义，不能仅凭字节数填写。
\end{solution}

\subsubsection*{章末习题}

\paragraph*{第四章章末第 1 题}
\begin{solution}
检查目标头文件目录、链接库、源码路径，以及配置时找到的 OpenCV 版本与安装位置。
\end{solution}

\paragraph*{第四章章末第 2 题}
\begin{solution}
共 600 字节；赋值共享数据，修改 b 的像素会影响 a，clone 才创建独立副本。
\end{solution}

\paragraph*{第四章章末第 3 题}
\begin{solution}
不能，应先判空并报告路径；imwrite 通常按文件扩展名选择编码器。
\end{solution}

\paragraph*{第四章章末第 4 题}
\begin{solution}
都可能是红色，需考虑色环两端；灰度值等于阈值时输出 0。
\end{solution}

\paragraph*{第四章章末第 5 题}
\begin{solution}
可先试中值滤波；开运算先腐蚀后膨胀，闭运算先膨胀后腐蚀。
\end{solution}

\paragraph*{第四章章末第 6 题}
\begin{solution}
轮廓坐标围成的连续矩形边长为 19 与 9，面积 171；200 是非零像素计数。
\end{solution}

\paragraph*{第四章章末第 7 题}
\begin{solution}
输出 50×30；绘图使用输出图像坐标，位置与尺度需匹配。
\end{solution}

\paragraph*{第四章章末第 8 题}
\begin{solution}
常用 Hamming 距离；还需对应关系和几何一致性检查，不能只看描述子距离。
\end{solution}

\paragraph*{第四章章末第 9 题}
\begin{solution}
相同通道数可能具有不同颜色顺序或含义；应按 encoding 转换并保留消息头，必要时捕获转换异常。
\end{solution}

\paragraph*{第四章章末第 10 题}
\begin{solution}
两个输出都变为 9；各变量指向同一底层像素数据。
\end{solution}

\section*{第五章：神经网络与深度学习}
\phantomsection
\addcontentsline{toc}{section}{第五章题解：神经网络与深度学习}
\subsubsection*{随堂练习}

\paragraph*{随堂练习 5.1}
\begin{solution}
令训练和验证的目标均为 [1,2,3]，两组预测均为 [2,3,4]，两组均方误差都是 1。训练误差只说明模型在见过的数据上的拟合，无法单独反映新样本表现；还需独立验证数据，并检查分布与划分是否合理。
\end{solution}

\paragraph*{随堂练习 5.2}
\begin{solution}
在所用环境执行：
\begin{tbash}
python -c "import sys, torch; print(sys.executable); print(torch.__version__); print(torch.ones(2,3).shape)"
\end{tbash}
记录实际解释器路径、版本和 \texttt{torch.Size([2, 3])}。\texttt{python -m pip} 明确使用该解释器的 pip，避免系统中另一个 pip 可执行文件把包安装到其它环境。
\end{solution}

\paragraph*{随堂练习 5.3}
\begin{solution}
空间尺寸依次为 $32\to28\to14\to10\to5$。最后为 16×5×5，每样本展平维度为 400；批量 2 时输出形状为 [2,400]。若接入完整 LeNet，第一层全连接的输入维度也必须由 256 改为 400，不能只改输入图片尺寸。
\end{solution}

\paragraph*{随堂练习 5.4}
\begin{solution}
在 random\_split 得到训练子集之后、创建 DataLoader 之前，对训练子集再取前 128 项：
\begin{tpython}
train_set = torch.utils.data.Subset(train_set, range(128))
\end{tpython}
训练循环改为一轮，batch\_size 保持 64，因此有两个训练批次。输入形状应为 [64,1,28,28]，输出为 [64,10]，标签为 [64]；损失应有限，反向传播和参数保存不报错。此步骤只验证流程，不用 128 个样本的准确率评价模型能力，确认后恢复完整训练集。
\end{solution}

\paragraph*{随堂练习 5.5}
\begin{solution}
原线性层权重形状为 [1,2]，新建 \texttt{nn.Linear(3,1)} 需要 [1,3]，load\_state\_dict 会报告 weight 尺寸不匹配。参数文件不能完整替代模型结构与预处理说明；应该恢复相同结构，再加载参数。单纯关闭严格键检查也不能让尺寸不同的权重变得兼容。
\end{solution}

\paragraph*{随堂练习 5.6}
\begin{solution}
把 x 改为 \texttt{torch.ones(32,4)}，保持计时方式、线程设置和设备不变。分别记录中位每批用时 $T_1,T_{32}$；每样本平均计算开销为 $T_1$ 和 $T_{32}/32$，吞吐率为 $1/T_1$ 和 $32/T_{32}$。这不等于单张图片实际响应延迟，因为批处理还可能等待凑满输入；不应预先假定具体加速倍数。
\end{solution}

\subsubsection*{章末习题}

\paragraph*{第五章章末第 1 题}
\begin{solution}
优先怀疑过拟合；减少样本不是通用修复，可检查数据划分、增强、正则化与早停。
\end{solution}

\paragraph*{第五章章末第 2 题}
\begin{solution}
24×24、12×12、8×8、4×4；16 个输出通道对应 256 个展平特征。
\end{solution}

\paragraph*{第五章章末第 3 题}
\begin{solution}
接收 logits；PyTorch 默认累积梯度，常规逐批更新需要先清零。
\end{solution}

\paragraph*{第五章章末第 4 题}
\begin{solution}
不能。eval 改变 Dropout、BatchNorm 等模块行为；no\_grad 关闭梯度记录，二者作用不同。
\end{solution}

\paragraph*{第五章章末第 5 题}
\begin{solution}
需要网络定义、类别顺序、输入形状、归一化和其它预处理约定；设备映射也需匹配。
\end{solution}

\paragraph*{第五章章末第 6 题}
\begin{solution}
吞吐率分别为 250 与约 667 样本/s；批量 8 吞吐更高，但每批等待和执行使单样本延迟不能直接这样比较。
\end{solution}

\paragraph*{第五章章末第 7 题}
\begin{solution}
批量 2 时输出为 2×400，因为最后特征为 16×5×5。
\end{solution}

\section*{第六章：CUDA 与 GPU 加速}
\phantomsection
\addcontentsline{toc}{section}{第六章题解：CUDA 与 GPU 加速}
\subsubsection*{随堂练习}

\paragraph*{随堂练习 6.1}
\begin{solution}
逐像素阈值的输出只依赖对应输入，容易分配给独立线程。累计前缀和的第 i 项依赖前面多个输入，朴素递推有跨元素依赖；可通过并行扫描算法加速，但不能简单让每个线程读取尚未计算完成的前一结果。
\end{solution}

\paragraph*{随堂练习 6.2}
\begin{solution}
能说明当前命令环境可查询 NVIDIA 驱动和设备，但不能说明 Toolkit 的编译工具已安装或在 PATH 中。应检查 Toolkit 安装位置与 nvcc 路径。也不能仅凭 nvidia-smi 正常断定应用的全部运行库、容器透传或 GPU 架构配置都正确。
\end{solution}

\paragraph*{随堂练习 6.3}
\begin{solution}
三个长度 4096 的 float 数组，共需 $3\times4096\times4=49152$ 字节设备内存，即 48 KiB。上传两个输入为 32768 字节，下载一个输出为 16384 字节，一次完整流程同样拷贝 49152 字节；不含运行时其它开销。
\end{solution}

\paragraph*{随堂练习 6.4}
\begin{solution}
块数为 $\lceil257/256\rceil=2$，共启动 512 个线程，其中 255 个线程的索引超出数据范围。保留 \texttt{if (i < n)}，运行带错误检查的完整示例，全部 257 个结果应为 3。只检查首元素不能发现尾部越界。
\end{solution}

\paragraph*{随堂练习 6.5}
\begin{solution}
每个算子各上传和下载一次，共六次传输，合计 $3+6\times2=15$ ms。中间结果留在设备端，仅统一上传与下载一次，总用时为 $3+2\times2=7$ ms。这里忽略重叠传输、启动与调度开销，节省来自减少数据搬运。
\end{solution}

\paragraph*{随堂练习 6.6}
\begin{solution}
无 GPU 时可以阅读源码、计算索引与字节数；若安装了 Toolkit，还可能完成编译，但这不代表程序能运行。GPU 枚举、kernel 执行、设备结果校验、同步错误检查和性能计时需要可访问的兼容 GPU 及驱动。检查记录应分“已执行”“缺设备”“缺工具”，不要把跳过写成通过。
\end{solution}

\subsubsection*{章末习题}

\paragraph*{第六章章末第 1 题}
\begin{solution}
逐像素阈值化更直接；递推存在数据依赖，需要重新设计并行算法。
\end{solution}

\paragraph*{第六章章末第 2 题}
\begin{solution}
不能。驱动可用与本机是否安装、配置 Toolkit 是两个问题。
\end{solution}

\paragraph*{第六章章末第 3 题}
\begin{solution}
上传 2×512×4，下载 512×4，总计 6144 字节。
\end{solution}

\paragraph*{第六章章末第 4 题}
\begin{solution}
需要 3 块，共 768 线程，多出 255 个线程。
\end{solution}

\paragraph*{第六章章末第 5 题}
\begin{solution}
同一 Block 内共享；部分线程不参与屏障可能违反同步要求，不能照此编写协作代码。
\end{solution}

\paragraph*{第六章章末第 6 题}
\begin{solution}
不一定，启动和执行错误不同；检查 cudaGetLastError，并通过同步等操作检查执行阶段错误。
\end{solution}

\paragraph*{第六章章末第 7 题}
\begin{solution}
前者说明驱动和设备状态及支持版本信息，后者说明编译工具版本；两者不同本身不等于错误，需核对兼容性。
\end{solution}

\section*{第七章：相机与成像原理}
\phantomsection
\addcontentsline{toc}{section}{第七章题解：相机与成像原理}
\subsubsection*{随堂练习}

\paragraph*{随堂练习 7.1}
\begin{solution}
记录设备完整型号、数据表版本或来源，以及最大分辨率、像元尺寸（通常为微米）、快门类型；例如记录格式可以是“型号：实际型号；分辨率：数据表值；像元尺寸：数据表值；快门：global/rolling/未知”。分辨率不等于物理感光面尺寸，同属 CMOS 也不能推断快门类型。缺失字段如实标未知。
\end{solution}

\paragraph*{随堂练习 7.2}
\begin{solution}
从 f/4 收到 f/5.6，进光量约减半。保持亮度和增益不变，曝光应从 1 ms 延长到约 2 ms；若按标称比值精算，为 $(5.6/4)^2\approx1.96$ ms。相同图像运动速度下，曝光内位移也近似加倍，运动模糊风险增加。
\end{solution}

\paragraph*{随堂练习 7.3}
\begin{solution}
泊松模型下标准差为 $\sqrt{900}=30$，相对噪声为 $30/900\approx3.33\%$。补光可以增加收到的光子数，改善此模型下的相对散粒噪声；数字增益只是同时放大已有信号与噪声，不能提供额外光子信息。
\end{solution}

\paragraph*{随堂练习 7.4}
\begin{solution}
像素数为 $1280\times720=921600$。Bayer8 为 921600 字节，BGR8 为 2764800 字节。文件还可能包含头部、元数据、行填充，或使用压缩编码，因此文件长度不一定等于原始像素区长度。
\end{solution}

\paragraph*{随堂练习 7.5}
\begin{solution}
当 $\gamma=1$，映射成为 $y=x$。对整数范围 0 到 255，使用本章构表公式并 round 后应逐项等于输入，可直接检查：
\begin{tpython}
gamma = 1.0
lut = [round(255*(i/255)**(1/gamma)) for i in range(256)]
assert lut == list(range(256))
\end{tpython}
\end{solution}

\paragraph*{随堂练习 7.6}
\begin{solution}
代入焦距 16 mm、光圈数 8、弥散圆 0.01 mm，得 $H=16^2/(8\times0.01)+16=3216$ mm。与光圈数 4 时的 6416 mm 比较，收小光圈使超焦距减小；在同焦距、同对焦距离与清晰度判据下，通常使景深增加。实际极小光圈还需考虑衍射。
\end{solution}

\paragraph*{随堂练习 7.7}
\begin{solution}
图像速度为 800 pixel/s，要求位移不超过 1 pixel，则曝光时间满足 $t\le1/800=0.00125$ s，即 1.25 ms。保持曝光时间和运动速度不变，仅提高增益不会减少位移；增益只可能在亮度允许时帮助选择更短的曝光。
\end{solution}

\paragraph*{随堂练习 7.8}
\begin{solution}
小角度与相同成像比例近似下，宽度为 $20\times12/8=30$ pixel。条件包括目标尺寸与姿态、距离、传感器像元尺寸、输出分辨率和裁剪缩放方式均不变；离轴畸变或较近拍摄时，简单比例可能不够准确。
\end{solution}

\paragraph*{随堂练习 7.9}
\begin{solution}
分别建立“主要正对”和“加入不同倾角、位置、边缘覆盖”的照片集合，保持分辨率、棋盘格内角点数和方格尺寸相同，分别运行标定并记录有效照片数、RMS、K 与畸变参数。正对照片约束不足可能得到不稳定参数，即使拟合这些照片的 RMS 很低。还应检查未参与标定的照片、参数合理性及重复标定的一致性，不能编造固定 RMS 数值。
\end{solution}

\paragraph*{随堂练习 7.10}
\begin{solution}
按题设坐标约定，新焦距为 $1000\times0.5=500$ pixel，主点为 $(640-80)\times0.5=280$ pixel。裁剪改变坐标原点，缩放改变像素尺度，所以只改 width/height 会导致投影位置错误。还需同步更新对应内参、投影矩阵和 ROI；若具体缩放库采用半像素中心约定，按其实际坐标映射处理。
\end{solution}

\paragraph*{随堂练习 7.11}
\begin{solution}
先确认原图发布和棋盘格规格，在标定命令的 --ros-args 之前加入 --no-service-check，采集多位置、多倾角数据，完成 CALIBRATE 后检查结果并点击 SAVE。复制归档、取得 YAML，按驱动文档配置标定文件路径，重启后检查 CameraInfo 与校正图。该选项只跳过服务检查，不会给驱动增加服务，也不会自动加载文件；若驱动不支持加载和发布参数，需要补充驱动实现或选用支持该功能的图像发布方案。未部署前不能把保存成功视为完整验收。
\end{solution}

\paragraph*{随堂练习 7.12}
\begin{solution}
$B=-(-96)/800=0.12$ m，即 12 cm。--approximate 只允许时间戳在指定容差内的图像组成一对，不会让曝光同时发生，也不补偿目标运动。动态实验应采用可靠的同步采集，或明确建模并验证时间补偿；不能靠放宽容差掩盖问题。
\end{solution}

\paragraph*{随堂练习 7.13}
\begin{solution}
错位为 $1000\times0.010=10$ pixel。重影来自多次曝光之间目标移动，融合时不同位置叠加；过曝是像素达到饱和而丢失亮部细节。缩短帧间隔或改善运动对齐可针对重影，降低曝光等措施针对过曝，不能混为同一问题。
\end{solution}

\subsubsection*{章末习题}

\paragraph*{第七章章末第 1 题}
\begin{solution}
光圈数为 4，即 f/4；分母使用入瞳直径。
\end{solution}

\paragraph*{第七章章末第 2 题}
\begin{solution}
约 3 pixel；在相同运动条件下缩短曝光可减少模糊。
\end{solution}

\paragraph*{第七章章末第 3 题}
\begin{solution}
都不能。全局快门统一曝光时序，增益改变信号放大，不改变入射光子数量。
\end{solution}

\paragraph*{第七章章末第 4 题}
\begin{solution}
分别为 10000 与 30000 字节；Bayer 每像素一个颜色采样，RGB 每像素三个通道。
\end{solution}

\paragraph*{第七章章末第 5 题}
\begin{solution}
不能。已经饱和而丢失的信息无法仅靠亮度映射恢复。
\end{solution}

\paragraph*{第七章章末第 6 题}
\begin{solution}
理想情况下像素内参基本不变，平移尺度变大 1000 倍；尺度由标定物尺寸决定。
\end{solution}

\paragraph*{第七章章末第 7 题}
\begin{solution}
0.012 m；同步误差会变成几何不一致。
\end{solution}

\paragraph*{第七章章末第 8 题}
\begin{solution}
0.0005 s，即 0.5 ms；真实相机需核对设置单位及是否受帧周期约束。
\end{solution}

\paragraph*{第七章章末第 9 题}
\begin{solution}
CALIBRATE 求解参数，SAVE 保存采样与参数文件，COMMIT 请求驱动通过 SetCameraInfo 接收结果。只有保存文件或服务调用成功还不够；重启驱动后应核对 CameraInfo 的尺寸、模型、焦距和畸变与保存结果一致，再观察匹配图像的校正效果。无服务时走手动文件加载路径。CameraInfo 的校正旋转 R 不能当作相机相对底盘的 TF 外参。
\end{solution}

\paragraph*{第七章章末第 10 题}
\begin{solution}
检查照片是否仅为相似正面姿态、是否漏掉边缘、角点是否模糊、标定板是否平整以及打印尺寸是否准确。保留不同姿态验证图，在固定内参和畸变的前提下估计该图位姿，再查看残差与校正直线；同时比较重复标定的参数稳定性。低训练 RMS 不保证视场覆盖和模型正确。校正后的图像应使用新内参或 P 中对应的内参部分及零畸变；若仍使用原图的 K、D，会把模型与图像混用，可能重复补偿畸变。
\end{solution}

\section*{第八章：信号处理与滤波}
\phantomsection
\addcontentsline{toc}{section}{第八章题解：信号处理与滤波}
\subsubsection*{随堂练习}

\paragraph*{随堂练习 8.1}
\begin{solution}
优先检查观测模型、零偏与标定，因为真实位置未改变而传感器整体偏移。持续 0.1 m 的偏差可能是系统性零偏，不能简单当作零均值随机噪声处理；可先用已知参考位置测量并校准，必要时把偏置纳入估计状态。
\end{solution}

\paragraph*{随堂练习 8.2}
\begin{solution}
序列变为 [9,11,9,15]，均值为 11，而最新真值为 15，输出误差为 -4。窗口使用了运动前的历史样本，虽让数值平滑，却滞后于新的真实状态。判断滤波好坏应同时比较真值误差与延迟。
\end{solution}

\paragraph*{随堂练习 8.3}
\begin{solution}
输入连续为 1、初值 0。系数 0.2 的前三步为 0.2、0.36、0.488；系数 1 的前三步为 1、1、1。小系数抑制突变更强但跟随更慢；系数 1 完全采用当前输入，不提供平滑。
\end{solution}

\paragraph*{随堂练习 8.4}
\begin{solution}
先验方差仍为 1.1，$K=1.1/(1.1+9)\approx0.108911$；更新估计 $x=0+2K\approx0.217822$，更新方差约为 0.980198。R 增大意味着观测更不确定，在本标量模型下对观测的权重减小。
\end{solution}

\paragraph*{随堂练习 8.5}
\begin{solution}
两个估计为 2 和 4，方差均为 1，权重各为 0.5，融合均值为 3；独立时融合方差为 0.5。同一帧派生的两个估计可能共享噪声，存在交叉协方差；忽略相关性会把同一份信息重复计算，使不确定性被低估。
\end{solution}

\paragraph*{随堂练习 8.6}
\begin{solution}
sigma 点的加权均值为 0。对 $y=2x+1$，加权输出均值为 $2\times0+1=1$，也等于 $f(E[x])$。该结论来自线性变换保持期望关系；任意非线性函数可能涉及更高阶统计量，有限 sigma 点不能保证全部精确表达。
\end{solution}

\subsubsection*{章末习题}

\paragraph*{第八章章末第 1 题}
\begin{solution}
前者归入过程模型不确定性，后者归入观测噪声；具体归类仍取决于定义的状态与模型。
\end{solution}

\paragraph*{第八章章末第 2 题}
\begin{solution}
依次为 0.5 和 0.875；系数越小通常越平滑也越滞后。
\end{solution}

\paragraph*{第八章章末第 3 题}
\begin{solution}
K=2/5=0.4，更新值为 4+0.4×5=6。
\end{solution}

\paragraph*{第八章章末第 4 题}
\begin{solution}
独立时为 2；完全相关时必须计入交叉协方差，不能直接套用独立公式。
\end{solution}

\paragraph*{第八章章末第 5 题}
\begin{solution}
EKF 在估计点用雅可比线性化，UKF 传播选定 sigma 点并重建统计量，二者都不是任意非线性问题的精确解。
\end{solution}

\paragraph*{第八章章末第 6 题}
\begin{solution}
创新量单位不一致，估计会严重偏离；应先统一单位与模型，不能用 R 掩盖尺度错误。
\end{solution}

\paragraph*{第八章章末第 7 题}
\begin{solution}
K=1.1/10.1，约 0.10891；x 约 0.21782，P 约 0.98020。测量权重减小。
\end{solution}

\section*{第九章：几何视觉与 Bundle Adjustment}
\phantomsection
\addcontentsline{toc}{section}{第九章题解：几何视觉与 Bundle Adjustment}
\subsubsection*{随堂练习}

\paragraph*{随堂练习 9.1}
\begin{solution}
深度由 2 m 变为 4 m，$u=800\times0.2/4+320=360$，$v=800\times0.1/4+240=260$。缩半的是相对主点的偏移，主点本身仍为 (320,240)，因此不能把原像素坐标 (400,280) 整体除以 2。
\end{solution}

\paragraph*{随堂练习 9.2}
\begin{solution}
先缩小到一半，$f_y=410$、$c_y=120$；再从顶部裁去 10 pixel，相当于图像原点下移 10，最终 $f_y=410$、$c_y=110$。这里按连续坐标缩放模型计算，实际重采样还需统一像素中心约定。
\end{solution}

\paragraph*{随堂练习 9.3}
\begin{solution}
列向量约定下取
\[
R=\begin{bmatrix}0&-1&0\\1&0&0\\0&0&1\end{bmatrix},\quad
p=\begin{bmatrix}0\\2\\3\end{bmatrix},\quad
t=\begin{bmatrix}1\\0\\0\end{bmatrix}.
\]
于是 $Rp+t=(-1,0,3)^T$。代入 $R^T(q-t)$ 应恢复原点坐标 (0,2,3)。
\end{solution}

\paragraph*{随堂练习 9.4}
\begin{solution}
保留原三维点顺序，只交换两项二维观测，然后重新求位姿并对所有点计算重投影 RMS。通常误差会明显增大，也可能返回不合理位姿或失败；具体值受求解器影响。成功返回只表示算法给出了结果，仍须核对点序、残差、正深度与位姿范围，不能把布尔返回值当成数据正确的证明。
\end{solution}

\paragraph*{随堂练习 9.5}
\begin{solution}
三个差值变为 2、2、8，因此 $t=(2+2+8)/3=4$。预测为 [4,5,6]，观测为 [2,3,10]，残差平方和为 $2^2+2^2+(-4)^2=24$。平方损失对大残差赋予很大代价，导致异常点把解从 2 拉到 4；可考虑异常点检测或适当的鲁棒损失。
\end{solution}

\paragraph*{随堂练习 9.6}
\begin{solution}
共有 3 个位姿顶点和 2 个路标顶点，即 5 个顶点；每个相机都观测两个点，共 6 条观测边。可以固定第一个相机位姿作为参考坐标。若是纯单目联合优化，还需处理整体尺度自由度；本题的图也不保证拥有足够的几何约束，固定参考并不自动保证所有参数可观。
\end{solution}

\paragraph*{随堂练习 9.7}
\begin{solution}
\texttt{numpy.linalg.qr(M)} 返回 Q 和上三角因子 U，应检查 \texttt{Q @ U} 重建 M。通常 Q 不等于原来的旋转 R，U 也不等于原内参 K，因为原模型是 K 在左、R 在右。可以打印两个差值矩阵或使用 allclose 验证，而不能只看到一个正交矩阵就称其为相机旋转。
\end{solution}

\paragraph*{随堂练习 9.8}
\begin{solution}
该矩阵本身是反射矩阵，未修正的 $UV^T$ 行列式为 -1。反射会改变坐标系手性，三维适当旋转必须满足正交性且行列式为 +1。按正文公式加入行列式修正后得到适当旋转；本例奇异值相同，最近旋转可能不唯一，不应要求某个固定的矩阵元素结果。
\end{solution}

\paragraph*{随堂练习 9.9}
\begin{solution}
例如令目标点为 $Q=PR^T+t$，其中 P 按行存放点，R 为绕 z 轴 90° 的旋转。只用质心差估计平移会得到 $(1/3,2,3)$，不能恢复原来的 $(1,2,3)$；三点残差平方和为 $8/3$。旋转改变点集相对质心的方向，单一平移无法补偿，需同时估计 R 和 t。
\end{solution}

\paragraph*{随堂练习 9.10}
\begin{solution}
角度改为 $-\pi/2$ 时，x 轴单位向量变为 $(0,-1,0)$，绕 z 轴的方向与正角相反。旋转矩阵的元素不能自由变化：需同时满足 $R^TR=I$ 和 $\det R=1$。逐元素相加通常破坏这些约束，应按统一约定用旋转复合或李代数扰动。
\end{solution}

\subsubsection*{章末习题}

\paragraph*{第九章章末第 1 题}
\begin{solution}
u=330 pixel；缩图后 fx=300、cx=150，未裁剪时对应 u=165。
\end{solution}

\paragraph*{第九章章末第 2 题}
\begin{solution}
p=Rᵀ(q-t)，因为旋转矩阵满足 R⁻¹=Rᵀ。
\end{solution}

\paragraph*{第九章章末第 3 题}
\begin{solution}
输入是三维点、对应二维点、内参与畸变；输出把物体系变到相机系的位姿。成功不代表对应正确，还需检查几何与残差。
\end{solution}

\paragraph*{第九章章末第 4 题}
\begin{solution}
3²+4²=25 pixel²；残差符号约定不影响平方和，但雅可比需保持一致。
\end{solution}

\paragraph*{第九章章末第 5 题}
\begin{solution}
分别代表相机姿态、三维点和图像约束；固定参考是为消除整体坐标变换的规范自由度。
\end{solution}

\paragraph*{第九章章末第 6 题}
\begin{solution}
使用 RQ；普通 QR 的相乘顺序是正交因子在左，不能直接交换。
\end{solution}

\paragraph*{第九章章末第 7 题}
\begin{solution}
使用 U diag(1,1,det(UVᵀ)) Vᵀ，要求奇异值按标准降序排列，以得到最近的适当旋转。
\end{solution}

\paragraph*{第九章章末第 8 题}
\begin{solution}
交替更新最近邻或其它对应关系，再估计刚体变换，直到满足收敛或停止条件。
\end{solution}

\paragraph*{第九章章末第 9 题}
\begin{solution}
结果为 (-1,0,0)；不能，左右扰动属于不同参考约定，雅可比必须匹配。
\end{solution}

\paragraph*{第九章章末第 10 题}
\begin{solution}
理想条件下旋转不变，平移约为原来的 1000 倍；这验证了位移单位由三维点单位决定。
\end{solution}

\section*{第十章：模型预测控制}
\phantomsection
\addcontentsline{toc}{section}{第十章题解：模型预测控制}
\subsubsection*{随堂练习}

\paragraph*{随堂练习 10.1}
\begin{solution}
改后响应为 $y_k=1-0.5^k$，第 3 步为 0.875，第 4 步为 0.9375，首次达到 0.9 需 4 步。原模型为 $y_k=1-0.8^k$，第 10 步约 0.8926、第 11 步约 0.9141，需 11 步。新模型更快接近目标，但这只是给定离散模型的比较。
\end{solution}

\paragraph*{随堂练习 10.2}
\begin{solution}
两频率分量的相对幅度被改变，高频分量仅保留原幅度的五分之一，叠加后的波形形状随之变化，属于幅频失真。如果还存在纯延时，整体会滞后，但纯延时本身不改变各频率分量之间的幅度比。
\end{solution}

\paragraph*{随堂练习 10.3}
\begin{solution}
宽度仍为 0.1 m，$\epsilon=\arctan(0.1/(2\times10))\approx0.005$ rad，约 $0.2865^\circ$，约为原来的一半。同样角度误差在更远处对应更大的横向偏移，所以更容易超出固定宽度目标的容限。
\end{solution}

\paragraph*{随堂练习 10.4}
\begin{solution}
代价为 $J(u)=(1-u)^2+9u^2$，求导得 $J'(u)=-2+20u$，无约束最优解为 0.1。它在 [-0.2,0.2] 范围内，因此约束解仍为 0.1，不触及限幅；控制惩罚增大使输入更温和。
\end{solution}

\paragraph*{随堂练习 10.5}
\begin{solution}
前三步控制均为 0.2，状态达到 0.6，再施加 -0.1 扰动后变为 0.5。重新规划的后三步控制为 0.2、0.15、0.075，状态依次为 0.7、0.85、0.925；若继续照抄无扰动时的后三步 0.2、0.1、0.05，状态为 0.7、0.8、0.85。闭环利用新状态修正未来输入，比盲目沿用旧序列更接近目标。
\end{solution}

\paragraph*{随堂练习 10.6}
\begin{solution}
日志至少包含当前周期编号、采集时间、开始/结束求解时间、返回状态、最大约束违反量与首个控制量。结果可用需满足求解器返回约定的成功状态、数值有限、约束容差满足要求、输入未过期且结果在控制截止时间内到达。任何一项失败都应进入已设计并验证的回退路径，不能只检查“函数正常返回”。
\end{solution}

\subsubsection*{章末习题}

\paragraph*{第十章章末第 1 题}
\begin{solution}
0.1 和 0.19；模型描述输入不能瞬间变成期望输出。
\end{solution}

\paragraph*{第十章章末第 2 题}
\begin{solution}
相位为 -0.1π rad，理想纯延时不改变幅值。
\end{solution}

\paragraph*{第十章章末第 3 题}
\begin{solution}
约为 0.08/(2×4)=0.01 rad；这只是几何近似，不是实车开火充分条件。
\end{solution}

\paragraph*{第十章章末第 4 题}
\begin{solution}
u|≤0.1 后呢？|无约束解为 0.25；约束解为 0.1。
\end{solution}

\paragraph*{第十章章末第 5 题}
\begin{solution}
现在执行 0.2；下一周期读取新状态重新求解，不盲目继续执行旧序列。
\end{solution}

\paragraph*{第十章章末第 6 题}
\begin{solution}
时域 0.1 s；求解超过周期，需检测超时与过期结果，并按已验证的回退策略处理。
\end{solution}

\paragraph*{第十章章末第 7 题}
\begin{solution}
控制仍应落在 [-0.2,0.2]；扰动后的控制根据新状态增大或维持限幅，而非直接沿用旧序列。
\end{solution}

\section*{第十一章：导航与 SLAM 初步}
\phantomsection
\addcontentsline{toc}{section}{第十一章题解：导航与 SLAM 初步}
\subsubsection*{随堂练习}

\paragraph*{随堂练习 11.1}
\begin{solution}
已有路径但没有速度时，先检查控制器是否激活、是否收到路径、局部代价地图与 TF 是否可用，以及进度/目标检查和下游速度接口。没有位姿时，先从定位输入、时间戳、TF 连通性和定位节点状态查起；没有坐标变换基础时，反复重算路径不能解决问题。
\end{solution}

\paragraph*{随堂练习 11.2}
\begin{solution}
停止原发布程序后，把 x 改为 2 再启动，tf2\_echo 应看到 x 平移为 2，y、z 为 0。不要让两个发布者同时发布同一个子坐标系。真实 map 到 odom 的变换用于随定位结果修正累计漂移，通常随时间变化；固定演示值不能长期替代定位输出。
\end{solution}

\paragraph*{随堂练习 11.3}
\begin{solution}
原例格心相对原点为 (0.525,1.025) m，加上新原点 (-2,-1)，得到 (-1.475,0.025) m。分辨率与栅格宽高共同决定地图覆盖的物理尺寸，例如宽 W 格时宽度为 $W\times resolution$。本题假定原点偏航为零，存在旋转时还需先旋转再平移。
\end{solution}

\paragraph*{随堂练习 11.4}
\begin{solution}
A 到 C 的直达代价改为 2，经 B 的代价仍为 $1+2=3$，所以最短距离为 2，走直接边。最少边只统计经过多少条边，最低代价则累加边权；原例直接边代价 5 时，一条边反而不如两条边的总代价 3。
\end{solution}

\paragraph*{随堂练习 11.5}
\begin{solution}
加入以下断言：
\begin{tpython}
assert fallback(["Failure", "Failure"]) == "Failure"
\end{tpython}
基础 Fallback 只传播子项状态；重试上限应由上层任务编排、重试装饰节点或明确的恢复策略决定，同时设超时与最终失败出口。不能由每个底层模块各自无限重试。
\end{solution}

\subsubsection*{章末习题}

\paragraph*{第十一章章末第 1 题}
\begin{solution}
先检查控制器状态、局部代价地图、TF、输入路径和进度条件，不能只反复调用全局规划。
\end{solution}

\paragraph*{第十一章章末第 2 题}
\begin{solution}
map 是全局参考，可发生定位修正；odom 局部连续但可能漂移；base\_link 固连机器人。
\end{solution}

\paragraph*{第十一章章末第 3 题}
\begin{solution}
格心为 (1.25,2.35) m；需区分栅格索引与图片行坐标。
\end{solution}

\paragraph*{第十一章章末第 4 题}
\begin{solution}
静态层提供已有地图，障碍层更新传感器障碍，膨胀层在障碍周围形成代价分布以反映通行余量。
\end{solution}

\paragraph*{第十一章章末第 5 题}
\begin{solution}
选择 A-B-C，总代价 4；最少边的直达路线并非最低代价。
\end{solution}

\paragraph*{第十一章章末第 6 题}
\begin{solution}
不执行，应返回 Running；失败才尝试后续分支。
\end{solution}

\paragraph*{第十一章章末第 7 题}
\begin{solution}
制动距离 v²/(2a)=2 m；还需加入感知、通信、控制延迟与余量，否则障碍被发现时可能已无法及时停车。
\end{solution}

\paragraph*{第十一章章末第 8 题}
\begin{solution}
到 C 的最短距离为 5，应走直接边。
\end{solution}

\endgroup
